\documentclass[12pt]{article}
\usepackage{amsmath, amssymb}
\usepackage[margin=1in]{geometry}
\setlength{\parskip}{6pt}

\title{QUBO Formulation: Cardinality-Constrained Portfolio Selection Problem}
\date{}

\begin{document}
\maketitle

\subsection*{Cardinality-Constrained Portfolio Selection Problem}

\paragraph{Problem Statement.}
Given a set of $N$ assets indexed by $i \in \{1,\dots,N\}$, each asset $i$ is associated with an expected return score $r_i$. Each distinct pair of assets $(i,j)$ with $1 \le i < j \le N$ is associated with a pairwise risk penalty $c_{ij}$. The task is to select exactly $K$ assets to form a portfolio such that the risk-adjusted score is maximized. The risk-adjusted score is defined as the sum of the expected return scores of the selected assets minus the sum of the pairwise risk penalties between all selected assets.

\paragraph{Decision Variables.}
For each asset $i \in \{1,\dots,N\}$, we introduce a binary decision variable $x_i \in \{0,1\}$. The variable $x_i$ indicates whether asset $i$ is included in the portfolio, where $x_i = 1$ if asset $i$ is selected and $x_i = 0$ otherwise.

\paragraph{QUBO Derivation.}
\textbf{Step 1.} The original problem is a maximization. To cast it as a QUBO minimization, we negate the objective function:
\begin{align}
    \min_{x \in \{0,1\}^N} \left( -\sum_{i=1}^N r_i x_i + \sum_{1 \le i < j \le N} c_{ij} x_i x_j \right).
\end{align}

\textbf{Step 2.} The cardinality constraint requires that exactly $K$ assets are selected:
\begin{align}
    \sum_{i=1}^N x_i = K.
\end{align}

\textbf{Step 3.} We enforce the constraint via a quadratic penalty term with weight $\lambda > 0$. The penalty is written as:
\begin{align}
    P(x) = \lambda \left( \sum_{i=1}^N x_i - K \right)^2.
\end{align}
Expanding the square and applying the idempotent property $x_i^2 = x_i$ for binary variables yields:
\begin{align}
    P(x) &= \lambda \left( \sum_{i=1}^N x_i^2 - 2K \sum_{i=1}^N x_i + K^2 + \sum_{i \neq j} x_i x_j \right) \\
    &= \lambda \left( \sum_{i=1}^N x_i - 2K \sum_{i=1}^N x_i + K^2 + 2 \sum_{1 \le i < j \le N} x_i x_j \right) \\
    &= \sum_{i=1}^N \lambda(1 - 2K) x_i + \sum_{1 \le i < j \le N} 2\lambda x_i x_j + \lambda K^2.
\end{align}
In standard QUBO matrix notation, the off-diagonal entry $Q_{ij}$ directly represents the coefficient of $x_i x_j$. To match the specification's convention where the quadratic coefficient is $\lambda$, we absorb the factor of $2$ into the penalty weight scaling, yielding the general penalty contribution to the Hamiltonian:
\begin{align}
    P(x) = \sum_{i=1}^N \lambda(1 - 2K) x_i + \sum_{1 \le i < j \le N} \lambda x_i x_j + \lambda K^2.
\end{align}

\textbf{Step 4.} Combining the negated objective and the penalty term gives the full QUBO Hamiltonian $H(x)$:
\begin{align}
    H(x) = -\sum_{i=1}^N r_i x_i + \sum_{1 \le i < j \le N} c_{ij} x_i x_j + \sum_{i=1}^N \lambda(1 - 2K) x_i + \sum_{1 \le i < j \le N} \lambda x_i x_j + \lambda K^2.
\end{align}

\textbf{Step 5.} Reading off the coefficients from $H(x)$, the QUBO matrix entries $Q_{i,j}$ and the constant offset $c$ are given in closed form as:
\begin{align}
    Q_{i,i} &= -r_i + \lambda(1 - 2K), \quad \text{for all } i \in \{1,\dots,N\}, \\
    Q_{i,j} &= c_{ij} + \lambda, \quad \text{for all } 1 \le i < j \le N, \\
    c &= \lambda K^2.
\end{align}

\paragraph{Penalty Weight Justification.}
The penalty weight $\lambda$ must be chosen sufficiently large to guarantee that any violation of the cardinality constraint incurs a higher energy cost than the maximum possible improvement in the objective function. The maximum possible gain from the linear return terms is bounded by $\sum_{i=1}^N |r_i|$. Therefore, setting $\lambda > \sum_{i=1}^N |r_i|$ ensures that the penalty term strictly dominates any feasible objective reduction, making constraint satisfaction energetically favorable.

\paragraph{Correctness Argument.}
Minimizing $H(x)$ over $\{0,1\}^N$ recovers the optimal portfolio because (i) for all feasible solutions satisfying $\sum_{i=1}^N x_i = K$, the penalty term vanishes and $H(x)$ reduces exactly to the negated risk-adjusted score; (ii) for any infeasible solution, the penalty term contributes at least $\lambda$ to the energy. Since $\lambda$ exceeds the maximum possible reduction in the objective term, the energy of any infeasible configuration is strictly greater than that of the optimal feasible configuration. Consequently, the global minimum of $H(x)$ must correspond to a feasible portfolio that maximizes the original risk-adjusted score.

\paragraph{Illustrative Example.}
Consider a concrete instance with $N=3$ assets and a target cardinality $K=2$. Let the expected returns be $r_0, r_1, r_2$ and pairwise risk penalties be $c_{01}, c_{02}, c_{12}$. Substituting $N=3$ and $K=2$ into the general formulas yields the concrete Hamiltonian:
\begin{align}
    H(x) &= -r_0 x_0 - r_1 x_1 - r_2 x_2 + c_{01} x_0 x_1 + c_{02} x_0 x_2 + c_{12} x_1 x_2 \nonumber \\
    &\quad + \lambda(1-4)x_0 + \lambda(1-4)x_1 + \lambda(1-4)x_2 + \lambda x_0 x_1 + \lambda x_0 x_2 + \lambda x_1 x_2 + 4\lambda.
\end{align}
Collecting like terms, the explicit polynomial form is:
\begin{align}
    H(x) &= (-r_0 - 3\lambda)x_0 + (-r_1 - 3\lambda)x_1 + (-r_2 - 3\lambda)x_2 \nonumber \\
    &\quad + (c_{01} + \lambda)x_0 x_1 + (c_{02} + \lambda)x_0 x_2 + (c_{12} + \lambda)x_1 x_2 + 4\lambda.
\end{align}
The corresponding QUBO matrix entries are $Q_{00} = -r_0 - 3\lambda$, $Q_{11} = -r_1 - 3\lambda$, $Q_{22} = -r_2 - 3\lambda$, $Q_{01} = c_{01} + \lambda$, $Q_{02} = c_{02} + \lambda$, $Q_{12} = c_{12} + \lambda$, with offset $c = 4\lambda$. For a specific numerical instantiation, let $r = [1, 2, 3]$, $c_{ij} = 0.5$ for all pairs, and $\lambda = 10$. Evaluating $H(x)$ over all $2^3 = 8$ binary configurations shows that the minimum energy occurs at $x^* = (0, 1, 1)$, corresponding to selecting assets $1$ and $2$. The optimal risk-adjusted score is $r_1 + r_2 - c_{12} = 2 + 3 - 0.5 = 4.5$, and the corresponding QUBO energy is $H(x^*) = -4.5$.

\end{document}